% !TeX root = ../../surface-lagrangian-torus-fibration.tex
% 2026-11-22 / v1.3: 固定 stage の有限 co-Gram 評価と最後の対角極限の分離．
\section{切り詰めた対数深度による全曲面の距離下界}
\label{ke:global-scalar-lower-section}

本節では $L=-\log|t|^2$，$R=\log L$，$\bar g_t=R^{-2}g_t$ とする．
Actual 全域 fibration の存在と component cap の評価は
定理\ref{ke:actual-global-fibration}および
補題\ref{ke:removed-corner-target-cap}を入力とする．
以下では clipped map を構成し，triple 外の単一 scalar 評価へ還元する．
その scalar 評価は節\ref{ke:sharp-normal-section}で証明した
normal Schwarz barrier によって供給される．さらに actual fibration の
基底値との全域比較を行い，同じ写像に対する全曲面の距離下界を得る．
外部 KE 研究の $\mathrm{(CP)}_d$，$\mathrm{(PBE)}_d$ を仮定しない．
以下の定数 $C_d$ は固定した次数と退化族に依存してよく，
$C_j$ はさらに固定 stage に依存してよい．

\subsection{全曲面上の clipped logarithmic map}

$\mathcal O_{\mathbf P^3}(1)$ に固定 Hermitian metric を取り，
各 $L_i$ を定数倍して
\[
 \ell_i=-\log\|L_i\|^2>1
\]
とする．$L_i=0$ では $\ell_i=+\infty$ と解釈する．
定数倍の積は defining equation の $F_d$ に吸収する．
$b_j$ を K2 の residual exhaustion scale，$m_{\rm old}=6b_j$ を旧 corner cut とする．
これとは独立に，固定 stage の全定数と core depth $N_j$ を決めた後で
lower clipping depth $m=\mathfrak m_j\gg m_{\rm old}+N_j+n_j+\mathfrak n_j$ と
plateau 定数 $K=K_j>2d$ を選ぶ．
いずれも stage とともに無限大へ発散させ，最後の slow diagonal に
$\log\mathfrak m_j+\log K_j=o(R_t)$ を課す．以下では $m=\mathfrak m_j$，$K=K_j$ と略記し，
\[
 M=\frac{L}{K},\qquad S=\log(M/m)=R-\log K-\log m,
 \qquad
 \Theta_i=\left[\frac{\log(\ell_i/m)}S\right]_0^1
\]
と置く．$[q]_0^1=\max(0,\min(q,1))$ である．

\begin{lemma}[Source 軸も含む大域的な cubical map]
\label{ke:global-clipped-depth-map}
十分小さい $|t|$ に対し
\[
 \Theta_t=(\Theta_1,\ldots,\Theta_d):X_t\longrightarrow\mathcal C_d
\]
は well-defined な continuous map で，固定 $t$ ごとに locally Lipschitz である．
少なくとも一つの coordinate は $1$ に等しく，非零 coordinate は高々三つである．
従って微分が非零となり得る active coordinates
\[
 A_t(x)=\{i:m<\ell_i(x)<M\}
\]
は高々二つである．Source 軸上でも $L_i=0$ の coordinate は plateau $1$ に入り，
その対数の発散による不連続性はない．
\end{lemma}

\begin{proof}
族の積方程式と誘導 Hermitian metric により，零でない点では
\[
 \sum_i\ell_i=L-\log\|F_d\|^2\ge L-C.
\]
従って $\max_i\ell_i\ge(L-C)/d>M$ である．
いずれかの $L_i$ が零ならこの結論は自明である．
一方，一般位置の任意の四平面には共通点がない．
各四元集合 $I$ について $\sum_{i\in I}\|L_i\|^2$ は
$\mathbf P^3$ 上で正の最小値を持つので，四番目に大きい $\ell_i$ は一様に有界である．
有限個の $I$ に共通の定数を越えて $m$ を取れば，$\ell_i>m$ は高々三つとなる．
Clip 関数は局所 Lipschitz であり，$L_i=0$ の開近傍では $\Theta_i=1$ である．
これで全主張が従う．
\end{proof}

\begin{lemma}[Co-Gram による全経路の長さ評価]
\label{ke:global-clipped-cogram-lower}
一つの $t$ と clipping parameters $m,K$ を固定し，$0<S\le R$ とする．
ある $\epsilon\ge0$ に対し，各非空集合 $I\subset\{1,\ldots,d\}$，
$|I|\le2$ の開集合
\[
 U_I=\{x\in X_t:m<\ell_i(x)<M\ \text{for all }i\in I\}
\]
上で，該当する covectors の co-Gram matrix がほとんど至る所で
\begin{equation}\label{ke:eq:global-active-upper}
 \left(\langle d\log\ell_i,d\log\ell_j\rangle_{g_t^{-1}}
       \right)_{i,j\in I}\le(1+\epsilon)\mathrm{Id}
\end{equation}
を満たすとする．$A_t(x)\ne\varnothing$ なら $I=A_t(x)$ と取ることで
全 active covectors を制御する．
このとき $X_t$ の任意の rectifiable curve $\gamma$ について
\[
 L_\square(\Theta_t\circ\gamma)
 \le\frac{R\sqrt{1+\epsilon}}S\,L_{\bar g_t}(\gamma).
\]
特に全点対で
\begin{equation}\label{ke:eq:global-raw-lower}
 R^{-1}d_{g_t}(x,y)
 \ge d_\square(\Theta_t x,\Theta_t y)
 -C_d\left(\epsilon+\frac{\log m+\log K}{R}\right).
\end{equation}
さらに $f_t=p_T\pi_t$ と
$\sup_xd_\square(f_t(x),\Theta_t(x))\le a_t$ なら，
右辺を $d_\square(f_t(x),f_t(y))$ に置き換えて誤差 $2a_t$ を加えられる．
これらは固定 parameter の不等式であり，$\epsilon$ の上限，
$\epsilon\to0$，又は $S/R\to1$ は仮定しない．
\end{lemma}

\begin{proof}
各 $U_I$ で co-Gram は smooth なので，ほとんど至る所の不等式は
連続性によりその全点へ延びる．他の coordinates が clipping level にある点も
$U_I$ から除いていないため，strict に active な残りの covectors に同じ評価が成立する．
従って ambient measure zero の集合に曲線が含まれる場合を除外していない．
Active region では $d\Theta_i=S^{-1}d\log\ell_i$ である．
Clip levels に沿って動く区間ではその coordinate の微分は零であり，
一般の absolutely continuous curve でもその level set 上でほとんど至る所同じことが成り立つ．
従って\eqref{ke:eq:global-active-upper}は
\[
 |(\Theta_t\circ\gamma)'|_{\ell^2}
 \le S^{-1}\sqrt{1+\epsilon}|\gamma'|_{g_t}
\]
を与える．有限 Euclidean cubical complex の curve の長さは，各 face の
Euclidean speed の積分である．これは relative interiors の逆像で計算し，
seams 上ではその法線 coordinate の微分がほとんど至る所零になることを用いればよい．
従って square seams や clipping levels の往復を含めて同じ積分評価が成り立つ．
全経路の長さについて下限を取ると，正確な有限評価
\[
 R^{-1}d_{g_t}(x,y)\ge q\,d_\square(\Theta_t x,\Theta_t y),
 \qquad q=\frac{S}{R\sqrt{1+\epsilon}}\in(0,1]
\]
を得る．$\eta=(\log m+\log K)/R=1-S/R\in[0,1)$ と置けば，
$\epsilon\ge0$ に対する $1-(1+\epsilon)^{-1/2}\le\epsilon/2$ より
\[
 1-q
 =\eta+(1-\eta)\{1-(1+\epsilon)^{-1/2}\}
 \le\eta+\epsilon/2.
\]
$D_d=\operatorname{diam}(\mathcal C_d,d_\square)<\infty$ を用いて
$(1-q)d_\square\le D_d(\eta+\epsilon/2)$ とすれば
\eqref{ke:eq:global-raw-lower}が従う．
最後に端点の三角不等式から
$d_\square(\Theta_t x,\Theta_t y)\ge d_\square(f_t(x),f_t(y))-2a_t$ であり，
$q\le1$ と合わせれば同じ誤差に $2a_t$ を加えるだけでよい．
領域間の往復回数を数えず，加法誤差も経路の各部分へ配分していない．
\end{proof}

\subsection{二つの active directions は既存 triple 比較で制御できる}

\begin{proposition}[大域下界の rank-one 評価への還元]
\label{ke:global-lower-rank-one-reduction}
$\mathcal T_{j,t}$ を，三つの異なる global depths が $b_j/2$ 以上となる領域とする．
この領域上では，既存 full-triple comparison と rank-one Schwarz estimate により
\eqref{ke:eq:global-active-upper}が
\[
 \epsilon_{j,t}=C(b_j^{-1}+K_j^{-2})
\]
で成立する．$X_t\setminus\mathcal T_{j,t}$ では active coordinate は高々一つである．
従って大域長さ下界のために新たに十分な解析入力は，次だけである：
\begin{equation}\label{ke:eq:single-active-normal-input}
 \sup_{\substack{x\notin\mathcal T_{j,t}\\\mathfrak m_j<\ell_i(x)<L/K_j}}
       |d\log\ell_i|_{g_t}(x)\le1+\varepsilon_t,
 \qquad \varepsilon_t\longrightarrow0.
\end{equation}
この入力は source を含む pair 領域の \emph{一方向の上 co-gradient bound} であり，
全 tensor の二側比較，angular derivatives，あるいは portal 上界を要求していない．
\end{proposition}

\begin{proof}
三つの $L_i$ が十分小さい点は，対応する一つの三重点の固定近傍に入る．
そこで exact triple coordinates $\zeta_i\zeta_j\zeta_k=t$ を取り，
$s_a=-\log|\zeta_a|^2$ とする．固定 smooth units と Hermitian factors の差により
\[
 \ell_a=s_a+h_a,
 \qquad |h_a|+|dh_a|_{g_t}\le C.
\]
この微分評価には，当該 triple 内での full-triple tensor 比較と
complete simplex model の $g_L^{\rm sf}\ge c g_{\rm FS}$ を使う．
ここで全曲面上の一様 projective lower bound を新たな入力にしていない．
従って active directions について
\[
 d\log\ell_a=d\log s_a+O_{g_t}(b_j^{-1}).
\]
ここで $|d\log s_a|_{g_t}\le1+O(b_j^{-1})$ は，
関連 KE 主稿 \texttt{simplex:lem:rank-one-schwarz} と
\texttt{coarse:thm:full-triple-tensor}から得る．

二つの active depths は $s_a,s_b\le L/K_j+O(1)$，
残る一つは $L-s_a-s_b$ である．
同稿 \texttt{simplex:prop:dominant-product} 又は
\texttt{simplex:cor:active-cogram} により，二方向の co-Gram の上界は
$\mathrm{Id}+O(b_j^{-1}+K_j^{-2})$ となる．
一方向だけの場合は上述の rank-one upper bound で十分で，
co-Gram が $1$ に近いという両側の結論まで要求する必要はない．
Smooth unit の誤差を加えて所要の上界を得る．

二つの active coordinates があれば，それらと plateau $1$ を持つ別の coordinate
の三つが全て $\mathfrak m_j\gg b_j$ より深い．従って必ず $\mathcal T_{j,t}$ に属する．
その補集合では一方向だけが残り，式\eqref{ke:eq:single-active-normal-input}で十分となる．
\end{proof}

\begin{corollary}[全域 co-Gram と actual 距離に関する Lipschitz 評価]
\label{ke:global-depth-lipschitz}
定理\ref{ke:sharp-global-one-active}の順序でパラメータを選ぶ．
このとき全 $X_t$ で，補題\ref{ke:global-clipped-cogram-lower}の仮定は
\[
 \epsilon_{j,t}\le C\bigl(\varepsilon_j+b_j^{-1}+K_j^{-2}\bigr)
\]
で成立する．特に任意の二点について
\begin{equation}\label{ke:eq:global-depth-lipschitz}
 d_\square(\Theta_t(x),\Theta_t(y))
 \le\frac{\sqrt{1+\epsilon_{j,t}}}{S}\,d_{g_t}(x,y).
\end{equation}
この主張は以下の fibration の値との比較を入力にしていない．
\end{corollary}

\begin{proof}
$\mathcal T_{j,t}$ 内では前命題を用い，外では active coordinate が高々一つ
なので定理\ref{ke:sharp-global-one-active}を用いる．境界では連続性と
principal submatrix によって同じ不等式を得る．
補題\ref{ke:global-clipped-cogram-lower}の全曲線評価で下限を取ればよい．
\end{proof}

\subsection{同じ大域 fibration の基底値との比較}

本項では命題\ref{ke:global-base-parameterization}で選んだ parameter 化を用い，
$f_{j,t}=p_T\pi_{j,t}$ とする．辺の向きを同節に合わせて
\[
 E_{ik}(u)=
 \begin{cases}e_i+u e_k,&0\le u\le1,\\
 (2-u)e_i+e_k,&1\le u\le2
 \end{cases}
\]
とする．Normal generator は $(1,-1)$ であり，component $i$ の側は
$\mu<0$，component $k$ の側は $\mu>0$ である．

\begin{lemma}[共通 profile と reciprocal moment の対数]
\label{ke:common-profile-log-depth}
系\ref{ke:common-neck-longitude}の共通 profile $\sigma_t$ について，
各固定 stage に対して十分大きい有限な $A_j$ を取ると，
$|\mu|\le\delta_j$ で
\[
\begin{array}{ll}
 \displaystyle
 \sigma_t(\mu)=\frac{\log(1/|\mu|)}R
       +O(\eta_{j,t}+R^{-1}),
       & \mu<-A_j/L,\\[4pt]
 \displaystyle
 \sigma_t(\mu)=2-\frac{\log(1/\mu)}R
       +O(\eta_{j,t}+R^{-1}),
       & \mu>A_j/L,\\[4pt]
 \displaystyle
 \sigma_t(\mu)=1+O_j(R^{-1}),
       &|\mu|\le A_j/L,
\end{array}
\]
となる．ここで $\eta_{j,t}=C(1+\log(N_j+1))/R$ であり，
最初の二式の定数は固定族と degree に依存してよい．
\end{lemma}

\begin{proof}
補題\ref{ke:common-edge-moment-profile}の $v$ を用いる．
$\mu\in[-\delta_j,\delta_j]$ に対応する model levels では
両 normal depths が $N_j$ と固定比以上に深く，
$N_j\gg m_{\rm old}$ である．
式\eqref{ke:eq:hexagon-edge-profile}と
式\eqref{ke:eq:edge-moment-weighted-derivative}から
\[
 \frac{d\nu_0}{dv}\asymp
 \frac1a+\frac1c\asymp\frac{1+e^{|v|}}L,
 \qquad \nu_0(0)=0.
\]
従って $|v|\ge1$ では $|\nu_0(v)|\asymp e^{|v|}/L$ であり，
$|\nu_0(v)|\le A_j/L$ では $|v|\le C_{A_j}$ である．
$u_{0,t}=1+v/R$ を用いれば表示式を得る．
Endpoints の rescaling は $O(\eta_{j,t})$ しか値を動かさず，
中央の固定区間では恒等なので最後の式も保たれる．
\end{proof}

この補題は共通 profile の独立な定量的説明である．以下の全点比較には，
normal 深度による中央帯の再分類に代えて，同一 moment の actual 接続を用いる．

\begin{lemma}[同一 moment を持つ旧 corner 辺への参照点]
\label{ke:same-moment-corner-reference}
固定 stage $j$ の各 pair $ik$ について，元の K2 neck の全残余領域，source discs，
および必要な全 end annuli を含む連結な固定 residual compact
$P_{ik,j}\Subset P^+_{ik,j}$ を，component cut $N$ より先に選べる．
Source discs を除いた部分も連結とする．ある有限な $N_{0,j}$ と
$C_j^{\rm neck}$ があり，$N\ge N_{0,j}$，$\delta=4\pi/N$ とし，その後 $t$ を
十分小さくすると次が成立する．

Pair $ik$ の構成後の neck の点 $x$ が actual normal moment
$\widehat\mu(x)=\mu\in[-\delta,\delta]$ を持つとする．任意の incident triple
$J=\{i,k,a\}$ について，旧 softmax 辺 $w_a=0$ 上に同じ moment を持つ全 torus
$\Lambda_{ik,J,t}(\mu)\subset A_{t,J}$ が存在し，その任意の点 $y$ に対して
\begin{equation}\label{ke:eq:same-moment-corner-reference}
 d_{g_t}(x,y)\le C_j^{\rm neck},\qquad
 P^{\rm old}_{t,J}(y)=E_{ik}(u_{J,t}(\mu))
\end{equation}
が成立する．定数 $C_j^{\rm neck}$ は $t,\mu,N$ に依存しない．
全 phase と $\mu=\pm\delta$ を含み，$x$ は source 又は saddle の特異 fiber 上に
あってよい．ここで接続経路は一つの fiber 又は moment level の中に留まる必要はない．
\end{lemma}

\begin{proof}
旧 cut を $m_{\rm old}=6b_j$ とする．
式\eqref{ke:eq:hexagon-edge-profile}より，旧辺の全 source parameter 範囲で
その residual 深度は
\[
 B=m_{\rm old}+r_a,\qquad
 0\le r_a=\max(r_i,r_k)/L<1
\]
である．従って，この全範囲を含む固定 end annuli と，marked tree 及び source
discs を含む元の neck の全残余領域を先に合わせて $P_{ik,j}$ を選べる．
End の moment-dependent な境界 graph の像も，K2 の固定 buffer によって
あらかじめ用意した広い residual domain の内部に入っている．少し広い
$P^+_{ik,j}$ も K2 で用意した residual domain 内に取れる．必要な annuli は
旧 trim $B=7b_j$ より内側にあり，$Z_{a,t}/Z_{a,0}\to1$ と固定 frame の差を
吸収する余裕を持たせる．この選択に $N$ 又は $[-\delta,\delta]$ は使わない．
ここで全旧辺を projected neck chart と同一視しているのではなく，その source
residual 範囲から annuli を先に固定している．

補題\ref{ke:common-edge-moment-profile}によれば，$N$ が十分大きいと
$|\mu|\le\delta$ に対応する旧辺の部分では両 normal 深度が $cN$ 以上となる．
$N\gg n_j$ と固定 frame の全定数を $N_{0,j}$ に含め，最後に $t$ を小さくすれば，
この必要な部分の全 tori は K2 の共通 plateau overlap に入る．そこでは
source residual coordinate が上で選んだ annulus を与える．
同補題の strict monotonicity が $\Lambda_{ik,J,t}(\mu)$ の存在と一意性を与え，
補題\ref{ke:corner-cap-disk}の proper arc の端点が $\mu=\pm\delta$ であることから，
両端を含めてこの torus は保持集合に属する．

系\ref{ke:neck-corner-saturated-union}と定理\ref{ke:staged-neck-averaging}により，
この全 torus 上の actual normal action と canonical moment は neck のものと
厳密に同じである．すなわち同じ normal Moser map $\Phi^N_t$ の前では
$\mu_N$，後では $\widehat\mu=\mu_N\circ(\Phi^N_t)^{-1}$ であり，
corner source では $\mu_\alpha=\mu_N\circ G$ と書けば，
式\eqref{ke:eq:staged-moment-identity}と K2 の flow 自然性により
\[
 \mu_\alpha\circ\psi_t=\mathcal A_2\mu_\alpha=\nu_J,
 \qquad \widehat\mu\circ C_{t,J}=\nu_J
\]
となる．第二段階の前後で同じ関数 $\mu_\alpha$ が点ごとに不変という意味ではなく，
この exact transport によって全平均 moment と actual moment を同定する．
新たな加法定数又は phase の
選択はない．相対葉族の補間も actual normal moment を保つので，構成後の $x$ と
$\Lambda_{ik,J,t}(\mu)$ は同じ actual moment 集合に入る．

固定基底に対する系\ref{ke:source-inclusive-same-moment-collapse}を適用する．
そこで用いる定理\ref{ke:fixed-base-neck-moser}と
命題\ref{ke:fixed-base-moment-connectors}は，固定基底の有限な下限以上の
normal cut $m$ と $m\le\vartheta M_t$ に対し，$m,t$ に独立な定数を持つ．
前者は固定 actual buffer と $C_P/m$ の短円評価，後者は固定した凸性の割線定数
$c_0,C_0,\vartheta,m_0$ による．従って固定 stage でも適用でき，この使用に
同じ stage 内で $m\to\infty$ とする仮定は要らない．必要な normal cut の下限は
K2 の $n_j$ の選択に含め，元の共通 moment interval を固定した後に
$[-\delta,\delta]$ をその内部へ入れる．Source からの全 moment 退出にも
この interval 全体に共通な定数がある．

以上の系は，$P_{ik,j}$ 上の同じ moment を持つ全点の合併に対して
$C+C_{P_{ik,j}}$ の ambient 直径上界を与える．同じ Moser map の両端の短い
trajectories の費用を加えれば，actual 像にも一つの $C_j^{\rm neck}$ を使える．
これらの定数は $N$ より先に固定され，後から moment interval を狭めても変わらない．
全 phase，異なる reduced 葉，source 固定点，saddle 葉も同じ集合評価に含まれる．
旧辺の projection 値は定義通り表示された $E_{ik}(u_{J,t}(\mu))$ である．
\end{proof}

以下では，この有限個の pair に対する参照点の定数と，
系\ref{ke:source-inclusive-same-moment-collapse}の同じ固定基底上の接続定数の
最大を $C_j^{\rm neck}$ とする．これを $N_j$ より先に固定し，component cut の
既存の下限に
\begin{equation}\label{ke:eq:neck-constant-depth-choice}
 C_j^{\rm neck}\le1+\log(N_j+1)
\end{equation}
を加える．これは許容される cut の選択条件であり，$C_j^{\rm neck}$ の
exhaustion に関する増大率を評価した結果ではない．後節の距離上界でも，
この同じ固定基底，定数，および cut 条件を用いる．

\begin{proposition}[Clipped depth と actual 基底値の全域比較]
\label{ke:global-depth-base-comparison}
式\eqref{ke:eq:neck-constant-depth-choice}を満たすよう $N_j$ を選ぶ．
Sharp normal 評価の下側 cut $\mathfrak m_j$ を，
固定 stage の core と変更 collar の非成分深度の全上界，
および保持 neck の残余 compact 上の非 pair 深度の全上界より
大きく取る．
この選択は定理\ref{ke:sharp-global-one-active}のパラメータ順序と両立する．
固定した次数・退化族と予め固定した chart・weights にのみ依存し，
$j,t$ に依存しない定数 $C$ に対して
\begin{equation}\label{ke:eq:global-depth-base-comparison}
 \sup_{x\in X_t}d_\square(\Theta_t(x),f_{j,t}(x))
 \le \frac{C\{1+\log\mathfrak m_j+\log K_j+\log(N_j+1)\}}R
\end{equation}
となる．$\Theta_t$ 自体が Lagrangian fibers 上で一定である必要はない．
\end{proposition}

\begin{proof}
\emph{Corner．}
旧 corner の source softmax は
\[
 s_a=m_{\rm old}+
     (L-3m_{\rm old})\frac{L^{w_a}}{\sum_bL^{w_b}}.
\]
$0\le w_a\le1$，$\max w_a=1$ より，その全域で
$|\log s_a-Rw_a|\le C\log(m_{\rm old}+1)$ である．
Global section norm との差は固定 smooth units なので有界である．
Dominant coordinate は旧写像と $\Theta_t$ の双方で $1$ となり，
他の coordinates は同じ三つの labels に限られる．
従って同じ square 内の線分を用いて，Moser--Hamiltonian 補正前の点 $y$ では
\[
 d_\square(\Theta_t(y),w)
 \le C\frac{\log\mathfrak m_j+\log K_j}R.
\]
ここまでの定数は有限個の固定 triple charts の unit bounds によるもので，
stage に依存しない．旧 corner map $C_{t,J}=G\Upsilon\Phi$ の変位にも
同じ一様性がある．実際，系\ref{ke:neck-corner-saturated-union}の証明では
入れ子領域の cutoffs を
$\log(b/b_j),\log(a/n_j),\log(c/n_j)$ の固定した滑らかな関数で作り，
その全固定階の模型 norm は一様である．相対化は moment interval に制限する前に
行われ，$\iota_j^{-1}$ による微分損失は生じない．
命題\ref{ke:quantitative-order-hamiltonian}と
定理\ref{ke:relative-corner-patch}の小摂動量を $\delta_{\rm av}$ と書くと，
$e=O(b_j^{-1})$，$\delta_{\rm av}=O(b_j^{-1})$ であり，
相対 Hamiltonian flow の軌跡長は模型計量で $O(b_j^{-2})$ となる．
全平均の Moser flow には定理\ref{ke:slow-cut-closed-corner-gh}の
生成場評価 $C/b_j$ を使えば十分である．また
命題\ref{ke:three-arm-coordinate-interpolation}の $G_s$ の軌跡長は
$O(e^{-cb_j}/b_j)$ である．全軌跡は最小深度が $c'b_j$ 以上の三重点域に留まり，
full-triple 比較はそこで actual 計量と模型計量の一様な二側比較を与える．従って
定理\ref{ke:slow-cut-closed-corner-gh}と同じ全 phase の領域で
\begin{equation}\label{ke:eq:corner-uniform-actual-displacement}
 \sup_{z\in\mathcal H_{t,m_{\rm old}}}
 d_{g_t}(z,C_{t,J}(z))\le C/b_j
\end{equation}
と評価でき，この $C$ も $j,t$ に依存しない．
補正前後の点を式\eqref{ke:eq:global-depth-lipschitz}で比較すると，
その追加誤差は $C/(b_jR)$ 以下である．この順序により，
補正の像で raw physical depth の微分を別に評価する必要はない．
命題\ref{ke:global-base-parameterization}を用いれば，
全保持 corner で所要の値比較が従う．

\emph{Core と変更 collar．}
補題\ref{ke:core-depth-upper}と固定 compact の graph convergence により，
component $i$ のこの領域では
\[
 \ell_a\le D_j\quad(a\ne i),\qquad \ell_i\ge L-D_j
\]
となる有限な $D_j$ がある．$\mathfrak m_j>D_j$，
$K_j>2$，その後 $t$ を十分小さく取れば $\Theta_t=e_i$ が厳密に成立する．
同じ基底 parameter 化の cap 評価より
$d_\square(f_{j,t},e_i)\le C\eta_{j,t}$ である．

\emph{同一 moment の参照点による neck の比較．}
共通 anchor interval 内の保持 neck の点 $x$ を取り，
$\mu=\widehat\mu(x)\in[-\delta_j,\delta_j]$ とする．
補題\ref{ke:same-moment-corner-reference}により，任意の incident triple $J$ の
旧 corner 辺の torus から参照点 $y$ を選べる．既証の
式\eqref{ke:eq:global-depth-lipschitz}と
$S=R-\log\mathfrak m_j-\log K_j\ge R/2$，および
式\eqref{ke:eq:neck-constant-depth-choice}により
\[
 d_\square(\Theta_t(x),\Theta_t(y))
 \le\frac{\sqrt{1+\epsilon_{j,t}}}{S}C_j^{\rm neck}
 \le\frac{C\{1+\log(N_j+1)\}}R .
\]
上の corner の議論は $y$ において
\[
 d_\square(\Theta_t(y),P^{\rm old}_{t,J}(y))
 \le C\frac{1+\log\mathfrak m_j+\log K_j}R
\]
を与える．さらに $P^{\rm old}_{t,J}(y)=E_{ik}(u_{J,t}(\mu))$ であり，
系\ref{ke:common-neck-longitude}の endpoint rescaling を含めた比較から
\[
 d_\square(P^{\rm old}_{t,J}(y),E_{ik}(\sigma_t(\mu)))\le C\eta_{j,t}.
\]
系\ref{ke:neck-base-value-control}を同じ actual moment $\mu$ に適用すると，
$d_\square(f_{j,t}(x),E_{ik}(\sigma_t(\mu)))\le C\eta_{j,t}$ である．
この四つの不等式を加えれば，neck 上で所要の比較が従う．Source 固定点，
pinched fiber の全点，branch fiber も同一 moment 集合の評価で扱っており，
それらの physical normal 深度を別に分類していない．

\emph{全被覆と依存の順序．}
命題\ref{ke:global-base-parameterization}と同じ共通 anchor で分割し，
core と変更 collar，保持 corner，anchor interval 内の保持 neck を用いる．
補題\ref{ke:anchor-trimmed-neck-domain}により，旧 neck の
$|\widehat\mu|>\delta_j$ の部分は対応する component の core 側にあり，
平行 link との差と補間支持は変更 collar に含められる．従って
$\sigma_t$ をその定義区間の外へ延長して用いる必要はない．
Corner との全交わりでは K2 の normal action と moment が同じである．
定理\ref{ke:actual-global-fibration}の全域被覆により，これらの領域で全 $X_t$ が
尽くされる．同一 moment の接続は
系\ref{ke:source-inclusive-same-moment-collapse}から直接用いており，
後節の全曲面距離上界又はその connector 予算は入力にしていない．
また $\Theta_t$ の Lipschitz 評価は本命題に先立って証明済みなので，
基底値比較と全域下界の間にも循環はない．
Core の深度上界 $D_j$ は lower cut と最後の $t$ 閾値にだけ用い，値誤差には
残らない．Neck の接続定数は既に $1+\log(N_j+1)$ に吸収したので，
有限個の固定 chart の定数を一つの $C$ にまとめれば結論を得る．
\end{proof}

\begin{theorem}[Actual fibration に対する全曲面の距離下界]
\label{ke:global-ambient-lower}
任意の固定 $d\ge5$ と本稿の一般性条件を満たす族を考える．
定理\ref{ke:actual-global-fibration}と
命題\ref{ke:global-base-parameterization}の同じ写像
$f_{j,t}=p_T\pi_{j,t}$ を用いる．
式\eqref{ke:eq:neck-constant-depth-choice}を含め，
固定 stage のパラメータを上述の順序で選べば，
全 $x,y\in X_t$ に対して
\begin{equation}\label{ke:eq:global-ambient-lower}
\begin{split}
 R^{-1}d_{g_t}(x,y)\ge {}&
 d_\square(f_{j,t}(x),f_{j,t}(y))
 -C_d\bigl(\varepsilon_j+b_j^{-1}+K_j^{-2}\bigr)\\
 &-\frac{C\{1+\log\mathfrak m_j+\log K_j+\log(N_j+1)\}}R .
\end{split}
\end{equation}
ここで $C_d,C$ は固定した次数・族と構成 data にのみ依存し，$j,t$ に依存しない．
従って $j=j(t)\to\infty$ を十分遅く選ぶと右辺の全誤差が消える．
左辺は $X_t$ 全体の任意の経路を許す ambient 距離であり，
protected-corner 距離ではない．
\end{theorem}

\begin{proof}
系\ref{ke:global-depth-lipschitz}で得た co-Gram 上界を
補題\ref{ke:global-clipped-cogram-lower}へ代入し，
命題\ref{ke:global-depth-base-comparison}で両端点を比較する．
全域の長さ評価を先に証明しているので，neck，corner，core 間の
経路の出入りの回数による加法誤差は生じない．
各固定 stage の全定数を決めた後，
\[
 R_t\ge j\{1+\log\mathfrak m_j+\log K_j+\log(N_j+1)\}
\]
を課すまで $t$ 閾値を小さくしてよい．これとは別に，各固定 compact の比較や
構成に必要な既存の全ての閾値も満たすようにする．
$\varepsilon_j\to0$，$b_j,K_j\to\infty$ と合わせれば最後の主張を得る．
この下界から fibration の全射性や距離上界を逆に推論してはいない．
表示された評価は許容された cut の下での誤差予算であり，
$\mathfrak m_j,K_j,N_j$ の有効な増大率，又は固定 compact の収束閾値は評価していない．
従って最適な収束速度を結論するものではない．
\end{proof}

\subsection{任意回の出入りを扱う protected-buffer 原理}

\begin{lemma}[加法誤差を持つ局所下界の正しい大域化]
\label{ke:buffered-local-lower-globalization}
$(X,d_X)$ を length space，$(Y,d_Y)$ を直径 $D$ の metric space とし，
$f:X\to Y$ を continuous map とする．$U_\alpha\subset V_\alpha$ を
$X=\bigcup_\alpha U_\alpha$ なる cover とし，$V_\alpha$ は開とする．
全 $V_\alpha$ 内の曲線 $c$ に対し
\[
 d_Y(f(c(0)),f(c(1)))\le\lambda L_X(c)+\epsilon,
 \qquad \lambda\ge1,
\]
かつ $U_\alpha$ から $X\setminus V_\alpha$ への first-exit curve の長さは
$b>0$ 以上と仮定する．このとき
\[
 d_X(x,y)\ge
 \frac{d_Y(fx,fy)-\epsilon}{\lambda+\epsilon/b}
 \ge d_Y(fx,fy)-D(\lambda-1+\epsilon/b)-\epsilon.
\]
従って局所誤差 $\epsilon\to0$ に加えて，
$\lambda\to1$ と $\epsilon/b\to0$ が十分である．
\end{lemma}

\begin{proof}
有限長の任意の曲線 $\gamma$ の出発点を含む $U_\alpha$ を選び，
最初に $V_\alpha$ を出る時刻までを一つの部分曲線とする．
その時刻の点を含む次の $U_\beta$ を選んで繰り返す．
完了した部分曲線はそれぞれ長さ $b$ 以上なので，その個数は
$L_X(\gamma)/b$ 以下である．無限回の first exits が有限長の中で起きることも
この下界により排除される．最後の一片を含めて局所不等式を加えると
\[
 d_Y(fx,fy)\le(\lambda+\epsilon/b)L_X(\gamma)+\epsilon.
\]
全経路で下限を取ればよい．元の曲線が二領域間を実際に何回往復したかは制限していない．
また $f(U_\alpha)$ と $f(\partial V_\alpha)$ の target separation が $b_0>\epsilon$
なら，局所不等式から $b=(b_0-\epsilon)/\lambda$ を得られる．
\end{proof}

既存の corner protection は，inner cut $m_{\rm old}$ と
outer cut $m_{\rm old}/2$ の差だけなら
正規化した幅が $O(R^{-1})$ である．一方，局所 distortion には
$(\log m_{\rm old}+\log R)/R$ が含まれる．従ってこのままでは
$\epsilon/b\to0$ は導けない．本節では，上記の scalar differential によって
全曲線に対する乗法的な長さ評価を先に得ることで，この問題を避けた．

\subsection{有限個の飽和した全被覆でも，加法誤差だけでは足りない}

\begin{example}[全被覆上の protected 比較が大域下界を与えない例]
\label{ke:full-cover-oscillating-bypass}
$Y=[0,1]\times[-1,1]$ に Euclidean metric を入れる．
$0<c<1$，整数 $N\to\infty$，$h_N=N^{-3}$ とし
\[
 a_\pm(x)=1\pm c\sin(2\pi Nx),\qquad
 A_\pm(x)=\int_0^x a_\pm(r)\,dr
\]
とする．$|A_\pm(x)-x|\le c/(\pi N)$ である．
$X_N$ は同じ rectangle に smooth metric
$g_N=dy^2+a_N(x,y)^2dx^2$ を入れたものとする．
$y\ge h_N$ では $a_N=a_+$，$y\le-h_N$ では $a_N=a_-$，
$|y|\le h_N/2$ では $a_N=1+c$ とし，間を単調に滑らかに補間する．
従って上半分では $a_N\ge a_+$，下半分では $a_N\ge a_-$ である．
$f_N:X_N\to Y$ は恒等写像とする．

$U_\pm=\{\pm y\ge0\}$ は $f_N$-飽和した二つの閉集合による全被覆で，
共通部分は一本の連結な interval である．全 fibers は一点なので径は零である．
上半分では
\[
 g_N\ge dA_+^2+dy^2.
\]
逆に transformed coordinates $(A_+(x),y)$ の線分を $y\ge h_N$ へ押し上げ，
両端に垂直線分を付ければ，追加費用は $2h_N$ 以下である．
従って全点対 $p,q\in U_+$ で
\[
 \big|d_{U_+}^{\rm int}(p,q)-d_Y(p,q)\big|
 \le C/N+2h_N.
\]
下半分も同様である．さらに protected regions
$V_+=\{y>-h_N/4\}$，$V_-=\{y<h_N/4\}$ 内でも，
対応する lower bound は同じ calibration $(A_\pm,y)$ によって成立する．
従って各 $U_\pm$ には誤差 $O(N^{-1})$ の protected distance 比較がある．

しかし $(0,0)$ から $(1,0)$ へ，$\sin(2\pi Nx)$ の符号ごとに
$y=-h_N$ と $y=h_N$ の安い側を使い，符号が変わるたびに垂直移動すると
\[
 d_{X_N}((0,0),(1,0))
 \le\int_0^1\min(a_+,a_-)\,dx+O(Nh_N)
 =1-\frac{2c}{\pi}+o(1).
\]
Target distance は $1$ なので大域下界の誤差は消えない．
この例では full cover，飽和性，connected overlap，zero fiber diameter，
および局所 protected distortion の全てが成立する．
不足するのは誤差を支配する protection であり，$\epsilon_N/b_N\asymp N^2$ である．
Actual KE family の反例ではなく，有限個の局所加法誤差を形式的に有限回だけ
足してよいとする推論への反例である．
\end{example}

\subsection{Scalar Chern--Lu の境界制御を省略できない理由}

局所正則 coordinate $z$ に対し，$s=-\log|z|^2$ と置くと
半平面への rank-one Chern--Lu は
$\Delta\log Q\ge Q-1$ という logistic inequality を与える．
Global section norm では $dd^c\ell_i$ が固定 projective curvature である．
当該 tube 上で $g_t\ge c g_{\rm FS}$ を別途確認できれば，その trace は有界となる．
ただしこの事実だけで\eqref{ke:eq:single-active-normal-input}は従わない．
微分不等式の誤差を $O(1/\ell_i)$ に整理することと，
その領域の全境界から最大原理を適用できることは別の条件である．

この区別は固定した smooth negative KE manifold の smooth divisor 近傍でも見える．
Metric が非退化な点の正則 defining coordinate を $z$ とすると
\[
 |d\log(-\log|z|^2)|_g
 \asymp\frac{1}{|z|\,|\log|z|^2|}\longrightarrow\infty.
\]
$dd^c(-\log|z|^2)=0$ であり，局所 Schwarz の微分不等式そのものも成立するが，
punctured domain の source metric は complete ではない．
最大原理の boundary control がないことが破綻の理由である．
これは本稿の退化族に対する否定例ではない．
節\ref{ke:sharp-normal-section}では，正則 projected normal frame を用いて
$\Delta s=0$ を保ち，source を含む比較領域を
$\mathfrak n_j\le s\le\kappa L$ として，その三種類の境界を全て評価した．
外側 $s=\mathfrak n_j$ の粗い有限上界 $A_j$ と内側 $s=\kappa L$ の上界 $C_j$ は，
\[
 B(s)=1+\varepsilon_j+A_j\mathfrak n_j/s+C_js/(\kappa L)
\]
という明示的な障壁により吸収される．残余境界では
full-triple の係数一の評価を用いる．
最終 active band を
$\mathfrak m_j\gg A_j\mathfrak n_j/\varepsilon_j$，
$s\le L/K_j$，$K_j\gg C_j/(\kappa\varepsilon_j)$ に取ることで
両境界の寄与が小さくなる．$s=\kappa L$ の比較境界と，
写像 $\Theta_t$ の plateau $L/K_j$ は区別している．
この境界評価と rank-two triple 比較を合わせたものが
系\ref{ke:global-depth-lipschitz}であり，局所 logistic inequality
だけから大域 Schwarz bound を推論したのではない．
